% =====================================================================
%  04-persistencia.tex - Aplicacion HTTP y persistencia transaccional
% =====================================================================

\section{Aplicación HTTP autenticada y persistencia transaccional}\label{sec:impl-http}

Esta sección reúne lo común a todos los dominios: transacción auditada, preparación de la base y límites de ejecución. Realiza el control de acceso de \hyperref[tab:cu-04]{CU-04} y concreta la petición protegida de la \cref{fig:peticion}. HTTP usa los puertos de flujo de la \cref{fig:hexagonal-completa}. Un procesador documental concentra la preparación criptográfica compartida con el flujo offline. El servidor compone almacén documental y auditoría de PostgreSQL con la TSA local, sin repositorio de archivos ni proveedor remoto como Cincel.

\subsection{Confirmación transaccional y continuidad de evidencia}\label{sec:impl-http-transacciones}

Una mutación auditada confirma su cambio y su evento de negocio en una sola transacción, o ninguno de los dos (\cref{fig:impl-operacion-auditada}).

\begin{figure}[!htbp]
  \centering
  \begin{tikzpicture}[
  tramo/.style={draw=ESCOMGris, dotted, line width=0.5pt},
  tramoL/.style={font=\diag\bfseries, text=ESCOMAzulOscuro, fill=white,
    anchor=west, inner sep=1.5pt}]
  \Actores{qad/1.4/{Qadra},api/4.6/{API (web)},app/7.9/{Servicio de aplicación},rd/11.1/{Redis},pg/14.4/{PostgreSQL}}{-11.1}

  % Preparación
  \draw[tramo] (0,-0.75) -- (16,-0.75);
  \node[tramoL] at (0.1,-0.75) {Preparación: sin escribir};
  \Msg{1.4}{4.6}{-1.3}{\textbf{1} mutación y token}
  \Paso{4.6}{-1.42}{2.7}{Admisión (\cref{fig:impl-admision})}
  \Msg{4.6}{7.9}{-2.2}{\textbf{2} caso de uso}
  \Msg{7.9}{11.1}{-2.6}{\textbf{3} valida la sesión}
  \Msg{7.9}{14.4}{-3.0}{cuenta activa, rol y acceso previo}
  \Paso{7.9}{-3.12}{3.0}{Forma canónica y SHA-256; según la operación, AES, RSA, TSA o ZIP}
  \Msg{7.9}{11.1}{-4.35}{\textbf{4} reautentica: misma identidad}

  % Transacción
  \draw[tramo] (0,-4.65) -- (16,-4.65);
  \node[tramoL] at (0.1,-4.65) {Transacción auditada};
  \Msg{7.9}{14.4}{-5.3}{\textbf{5} READ COMMITTED y bloqueo\\ consultivo común}
  \Paso{14.4}{-5.42}{2.75}{Revalida cuenta, rol, pertenencia, estado del expediente y revisión esperada}
  \Paso[pasoOsc]{14.4}{-6.72}{2.75}{Escribe cambio, evento encadenado y, si aplica, recibo; confirma juntos}
  \Paso[pasoError]{7.9}{-5.42}{3.0}{Rechazo en 3 a 5 (401, 403, 404, 409, 400 o 422): la transacción no se abre o se revierte; no quedan filas ni evento de negocio}
  \Msg[msgR]{14.4}{7.9}{-7.9}{confirmado}
  \Msg[msgR]{7.9}{4.6}{-8.28}{revisión y, si existe, recibo}

  % Después de confirmar
  \draw[tramo] (0,-8.58) -- (16,-8.58);
  \node[tramoL] at (0.1,-8.58) {Después de confirmar};
  \Msg{4.6}{14.4}{-9.08}{\textbf{6} evento de acceso en una inserción separada}
  \Paso{14.4}{-9.2}{2.75}{Concedido; en un rechazo, denegado o rechazado}
  \Paso[pasoError]{4.6}{-9.2}{2.7}{Si esa inserción falla: error genérico 500}
  \Msg[msgR]{4.6}{1.4}{-10.2}{\textbf{7} respuesta}
  \Msg{1.4}{4.6}{-10.85}{\textbf{8} sin estado o 5xx:\\ consulta, no reenvía}
\end{tikzpicture}

  \caption{Recorrido de una mutación auditada. AES, RSA, TSA y ZIP corren antes de la transacción; el servicio vuelve a autenticar con la misma identidad y la transacción (paso 5) revalida el acceso y que el documento sea del expediente, y confirma juntos cambio, evento y, si existe, recibo. Un rechazo en los pasos 3 a 5 no deja filas; ante una respuesta incierta, Qadra consulta el resultado exacto en lugar de reenviar.}
  \label{fig:impl-operacion-auditada}
\end{figure}

El bloqueo consultivo común se toma antes de leer cabezas o bloquear filas, lo que serializa confirmaciones y cabeza de auditoría; secuencia y unicidad administrativa rechazan un aislamiento incompatible. La misma regla cubre la creación de expediente con su asignación inicial, los cambios de pertenencia y las mutaciones de identidad; listado y detalle confirman su evento de lectura en la transacción de autorización. Las operaciones que compiten se verifican en la \cref{sec:pruebas-transacciones}.

\subsection{Preparación administrativa e importación}\label{sec:impl-http-migracion}

El esquema durable se organiza en familias de migraciones numeradas de 0001 a 0048 (\cref{tab:anexo-migraciones}) y se prepara con el comando de migración sobre una conexión administrativa (\cref{sec:cli-database}). El servidor no ejecuta DDL y no arranca si su rol operativo puede administrar o reescribir auditoría, documentos, recibos o la protección de evidencia; ese rol lee y anexa eventos sin actualizarlos, eliminarlos ni truncarlos. Administradores y propietarios de PostgreSQL permanecen en la base de confianza.

El comando de importación migra el repositorio local anterior (\cref{fig:impl-http-importacion}; precondiciones y procedimiento en la \cref{sec:cli-database}). La primera importación exige vacías las tablas de documentos, clasificación, directorio, revisiones administrativas, registros iniciales y auditoría, y expedientes destino preparados administrativamente como expedientes de línea base, con hechos de creación conocidos y marcador de revisión inicial nulo; crearlos por HTTP y borrar revisiones o eventos no es válido. Antes de insertar se instalan barreras en el directorio de documentos y en la bitácora JSON Lines del origen, y tras confirmar, marcadores de migración; el respaldo operativo incluye PostgreSQL completo y protege por separado KEK y material criptográfico.

\begin{figure}[!htbp]
  \centering
  \begin{tikzpicture}[
  est/.style={caja, text width=3.0cm, minimum height=1.25cm},
  fin/.style={cajaOsc, text width=3.0cm, minimum height=1.25cm},
  fallo/.style={cajaError, text width=3.0cm, minimum height=1.1cm}]

  % Camino principal
  \node[est] (s0) at (1.6,0) {\textbf{Origen detenido}\\ respaldo y mapa completo de documento a expediente};
  \node[est] (s1) at (5.75,0) {\textbf{Inspección}\\ valida identidad, AAD, digest, firma, sello y cadena sin escribir};
  \node[est] (s2) at (9.9,0) {\textbf{Barreras}\\ migración pendiente en documentos y bitácora del origen};
  \node[fin] (s3) at (14.05,0) {\textbf{Confirmada}\\ documentos, prefijo histórico, evento y recibo en una transacción};
  \draw[flecha] (s0) -- (s1);
  \draw[flecha] (s1) -- node[etq, above=1pt]{aplicar} (s2);
  \draw[flecha] (s2) -- node[etq, above=1pt]{commit} (s3);

  % Fallos
  \node[fallo] (e1) at (5.75,-2.4) {Rechazo: el origen se conserva para resolución manual};
  \node[fallo] (e2) at (9.9,-2.4) {Fallo antes del commit: las barreras mantienen el origen bloqueado};
  \node[fallo] (e3) at (14.05,-2.4) {Fallo de marcadores: confirmada, con las fuentes aún bloqueadas};
  \draw[flechaError] (s1) -- (e1);
  \draw[flechaError] (s2) -- (e2);
  \draw[flechaError] (s3) -- (e3);

  % Cierre
  \node[fin] (s4) at (14.05,-4.75) {\textbf{Origen migrado}\\ marcadores de migración instalados};
  \node[est, text width=6.6cm] (s5) at (7.85,-4.75) {\textbf{Arranque del servidor sobre el origen preservado}: reconcilia marcadores, recibo, asociaciones, documentos y cadena completa; rechaza una restauración parcial};
  \draw[flecha] (e3) -- node[etq, left=2pt, align=right]{repetir mismas fuentes\\ y mapa: reconcilia\\ sin duplicar} (s4);
  \draw[flecha] (s3.east) -- ++(0.3,0) |- (s4.east);
  \draw[flecha] (s4) -- (s5);
\end{tikzpicture}

  \caption{Estados de una importación: las barreras se instalan antes de insertar y los marcadores después de confirmar; cualquier fallo deja el origen bloqueado, ningún rechazo reemite firmas ni sellos y repetir con las mismas fuentes y mapa completa sin duplicar.}
  \label{fig:impl-http-importacion}
\end{figure}

Las migraciones de dominio preservan la evidencia ya registrada (\cref{tab:impl-http-migraciones-dominio}). Las funciones de canon y secuencia resuelven sus referencias contra el esquema de instalación, así que una restauración ejecuta sus comprobaciones sin depender de la ruta de búsqueda de la sesión y sin eliminar restricciones.

\begin{table}[t]
  \centering\small
  \caption{Migraciones de los dominios documental y administrativo.}
  \label{tab:impl-http-migraciones-dominio}
  \begin{tabularx}{\linewidth}{@{}L{2.35cm} X X@{}}
    \toprule
    \textbf{Migración} & \textbf{Añade sin alterar evidencia} & \textbf{Arranque y respaldo} \\
    \midrule
    Versiones (\ref{sec:impl-http-versiones}) & Claves y raíces, sin recifrar ni tocar firma, sello, certificados, CRL ni prefijo; repetible; reconcilia por UUID y número originales & Claves, triggers, permisos, primera versión e inventario; respaldo con raíces y todas las versiones; ningún servidor que suponga UUID único \\
    Clasificación (\ref{sec:impl-http-clasificacion}) & Exige UTF-8; tabla y funciones; lo existente queda en revisión cero; el rol operativo solo lee e inserta revisiones, sin poseer sus protecciones & Restricciones, triggers, privilegios y continuidad \\
    Administrativa (\ref{sec:impl-http-administracion}) & Revisiones, sin convertir referencias en NUC ni fabricar etapas; altas con R1 por clave foránea diferida; el rol operativo no inserta el marcador de línea base & Canon, huella, continuidad, perfil irreversible y etapa inicial; respaldo con filas administrativas y registro inicial \\
    \bottomrule
  \end{tabularx}
\end{table}

Un recibo de importación confirmado admite revisiones administrativas posteriores sin duplicar eventos ni alterar el prefijo importado. Un ensayo automatizado cubre repositorio anterior, importación y restauración (\cref{sec:pruebas-migracion}).

\subsection{Límites de ejecución y verificación integrada}\label{sec:impl-http-limites}

Los adaptadores PostgreSQL y Redis son síncronos, así que HTTP ejecuta los casos de uso en trabajadores bloqueantes de Tokio con presupuesto compartido (\cref{fig:impl-admision}); los cupos se ajustan con las opciones de la \cref{sec:cli-serve}.

\begin{figure}[t]
  \centering
  \begin{tikzpicture}[
  et/.style={caja, text width=3.2cm, minimum height=1.25cm},
  alt/.style={caja, text width=3.2cm, minimum height=1.1cm}]

  \node[et] (a) at (1.8,0) {\textbf{Cupo de peticiones}\\ 8 por defecto; se toma antes de leer el cuerpo};
  \node[et] (b) at (5.9,0) {\textbf{Cuerpo acotado}\\ límite propio de cada familia de rutas};
  \node[et] (c) at (10.0,0) {\textbf{Espera de trabajador}\\ 2 por defecto; conserva el cupo de petición};
  \node[et, cajaOsc] (d) at (14.1,0) {\textbf{Caso de uso síncrono}\\ PostgreSQL y Redis en un hilo bloqueante};
  \draw[flecha] (a) -- (b);
  \draw[flecha] (b) -- (c);
  \draw[flecha] (c) -- (d);

  \node[alt, cajaError] (ea) at (1.8,-2.2) {Sin cupo: rechazo inmediato (503), sin cola ni ejecución};
  \node[alt, cajaGris] (ec) at (10.0,-2.2) {Cancelación en espera: se retira sin ejecutar};
  \node[alt, cajaGris] (ed) at (14.1,-2.2) {Cancelación en ejecución: el permiso sigue tomado hasta terminar; lo confirmado no se revierte};
  \draw[flechaError] (a) -- (ea);
  \draw[flechaSuave] (c) -- (ec);
  \draw[flechaSuave] (d) -- (ed);

  \node[caja, cajaClara, text width=7.3cm, minimum height=0.9cm] (h) at (3.85,-3.85)
    {La consulta de salud queda fuera de la admisión. El contenido documental verificado toma además un permiso de entrega hasta liberar sus bytes; sin él, también \texttt{503}};
  \node[caja, cajaClara, text width=7.3cm, minimum height=0.9cm] (s) at (12.1,-3.85)
    {Ante una señal detiene HTTP y los consumidores: deja de admitir y retiene el trabajo bloqueante iniciado; no hay plazo máximo si una dependencia síncrona no responde};
\end{tikzpicture}

  \caption{Admisión y límites de ejecución. Solo agotar un cupo (de petición o de entrega) produce \texttt{503} inmediato; la espera de trabajador queda acotada por el cupo de petición.}
  \label{fig:impl-admision}
\end{figure}

Cada familia de rutas limita su cuerpo JSON en su propia sección; identidad y expedientes básicos usan 16~KiB. Los errores responden 401 por credenciales o sesión, 403 por permisos, 404 por ausencia o aislamiento, 409 por conflicto, 400 o 422 por entrada inválida, 429 por límite de inicio de sesión y 500 por fallo interno, sin secretos. Ensayos con servicios desechables ejercitan este flujo (\cref{sec:pruebas-transacciones,sec:pruebas-backends}).

\paragraph{Límites.} La sesión en Redis y la transacción PostgreSQL son pasos separados: no se afirma atomicidad distribuida. El arranque no autentica definiciones SQL sustituidas por un administrador ni detecta en general la eliminación de toda la historia o una restauración antigua coherente. Los servicios escuchan solo en la interfaz de retorno local; la agrupación de conexiones, la evaluación de capacidad, los plazos operativos de todas las dependencias y el anclaje externo de la auditoría conservan aceptación pendiente.

\Fuentes{\rutarepo{docs/http-api.md}, \rutarepo{docs/database-operations.md} y los ADR 0015 y 0016.}
